% GENERATED by scripts/build-lecture-slides.py; do not edit. Portable profile: XeTeX, theme and fonts from the core.
\makeatletter\IfFileExists{core/beamer/beamerthememx.sty}{\def\mxcorepath{core/}\def\input@path{{core/beamer/}{core/}}}{\def\mxcorepath{../core/}\def\input@path{{../core/beamer/}{../core/}}}\makeatother
\documentclass[aspectratio=169,10pt,t]{beamer}
\usetheme{mx}
\usepackage{graphicx,booktabs}
\usepackage[french]{babel}
\setlength{\parskip}{.45em}
\hypersetup{pdfauthor={},pdfsubject={Lecture slides},pdfkeywords={},pdftitle={Barycentres et lignes de niveau}}
\providecommand{\N}{\mathbb{N}}
\providecommand{\Z}{\mathbb{Z}}
\providecommand{\D}{\mathbb{D}}
\providecommand{\Q}{\mathbb{Q}}
\providecommand{\R}{\mathbb{R}}
\providecommand{\C}{\mathbb{C}}
\providecommand{\vect}[1]{\overrightarrow{#1}}
\providecommand{\vi}{\vec{\imath}}
\providecommand{\vj}{\vec{\jmath}}
\providecommand{\vk}{\vec{k}}
\providecommand{\norm}[1]{\left\lVert #1\right\rVert}
\providecommand{\vc}[2]{\begin{pmatrix}#1\\#2\end{pmatrix}}
\providecommand{\deter}[2]{\det\left(#1,#2\right)}
\providecommand{\abs}[1]{\left\lvert #1\right\rvert}
\providecommand{\dd}{\,\mathrm{d}}
\providecommand{\e}{\mathrm{e}}
\providecommand{\Card}{\operatorname{Card}}
\providecommand{\E}{\operatorname{E}}
\providecommand{\V}{\operatorname{V}}
\providecommand{\degre}{^{\circ}}
\providecommand{\un}[1]{\,\mathrm{#1}}
\providecommand{\sci}[2]{#1\times10^{#2}}
\providecommand{\pgcd}{\operatorname{PGCD}}
\providecommand{\ppcm}{\operatorname{PPCM}}
\title{Barycentres et lignes de niveau}
\subtitle{Cours et exemples}
\date{}
\author{}
\begin{document}
\mxtitleframe
\begin{frame}{Objectifs}
Construire et calculer un barycentre, regrouper des points pondérés et déterminer des lieux géométriques.
\end{frame}
\begin{frame}{Deux, trois ou quatre points pondérés}
Pour des points \(A_i\) et des poids réels \(\alpha_i\) dont la somme \(s\) est non nulle, le barycentre \(G\) est l’unique point tel que \(\sum_i\alpha_i\overrightarrow{GA_i}=\vec0\).
\end{frame}
\begin{frame}{Deux, trois ou quatre points pondérés}
Pour tout point \(O\), \(\overrightarrow{OG}=\dfrac1s\sum_i\alpha_i\overrightarrow{OA_i}\). Pour tout \(M\),
\[
\sum_i\alpha_i\overrightarrow{MA_i}=s\overrightarrow{MG}.
\]
Les coordonnées de \(G\) sont les moyennes pondérées des coordonnées des points. Multiplier tous les poids par un même réel non nul ne change pas \(G\).
\end{frame}
\begin{frame}{Deux, trois ou quatre points pondérés}
Par Chasles, \(\overrightarrow{GA_i}=\overrightarrow{GO}+\overrightarrow{OA_i}\). L’équation du barycentre impose \(s\overrightarrow{OG}=\sum\alpha_i\overrightarrow{OA_i}\) : un unique vecteur définit un unique point. La réduction résulte de \(\overrightarrow{MA_i}=\overrightarrow{MG}+\overrightarrow{GA_i}\).
\end{frame}
\begin{frame}{Deux, trois ou quatre points pondérés}
Pour \((A,1),(B,2)\), \(\overrightarrow{AG}=\tfrac23\overrightarrow{AB}\). Avec des poids \(1,-2\), la somme vaut \(-1\) et \(\overrightarrow{AG}=2\overrightarrow{AB}\) : le barycentre peut être à l’extérieur du segment. Si les poids sont \(1,-1\) et \(A\ne B\), aucun point ne satisfait la définition, car la somme vectorielle est \(\overrightarrow{BA}\ne\vec0\).
\end{frame}
\begin{frame}{Deux, trois ou quatre points pondérés}
Un groupe de points dont la somme des poids est non nulle peut être remplacé par son barycentre affecté de cette somme. La somme totale doit rester non nulle.
\end{frame}
\begin{frame}{Deux, trois ou quatre points pondérés}
La réduction de ce groupe remplace exactement sa contribution dans la somme des vecteurs. L’équation finale ne change donc pas.
\end{frame}
\begin{frame}{Sommes de distances au carré}
Pour le barycentre \(G\) et \(s\ne0\),
\[
\sum_i\alpha_i MA_i^2=sMG^2+\sum_i\alpha_i GA_i^2.
\]
Les poids peuvent être négatifs; on ne doit alors pas supposer que la somme est positive.
\end{frame}
\begin{frame}{Sommes de distances au carré}
Développer le carré de \(\overrightarrow{MA_i}=\overrightarrow{MG}+\overrightarrow{GA_i}\), multiplier par \(\alpha_i\) et sommer. Le terme croisé vaut \(2\overrightarrow{MG}\cdot\sum\alpha_i\overrightarrow{GA_i}=0\).
\end{frame}
\begin{frame}{Reconnaître les lieux}
\(MA=k\) donne un cercle si \(k>0\), le point \(A\) si \(k=0\), l’ensemble vide si \(k<0\).
Pour \(\overrightarrow{v}\ne\vec0\), \(\overrightarrow{MA}\cdot\overrightarrow{v}=0\) est la droite par \(A\) perpendiculaire à \(\overrightarrow{v}\); pour \(\overrightarrow{v}=\vec0\), tout le plan.
Si \(I\) est milieu de \([AB]\),
\[
\overrightarrow{MA}\cdot\overrightarrow{MB}=MI^2-\dfrac{AB^2}{4}.
\]
Le lieu \(\overrightarrow{MA}\cdot\overrightarrow{MB}=k\) est donc un cercle, un point ou vide selon le signe de \(k+AB^2/4\).
\end{frame}
\begin{frame}{Reconnaître les lieux}
Écrire \(\overrightarrow{MA}=\overrightarrow{MI}+\overrightarrow{IA}\) et \(\overrightarrow{MB}=\overrightarrow{MI}-\overrightarrow{IA}\). Le produit vaut \(MI^2-IA^2\), et \(IA=AB/2\).
\end{frame}
\begin{frame}{Reconnaître les lieux}
Supposons \(A\ne B\) et cherchons \(MA=kMB\). Pour \(k<0\), aucun point; pour \(k=0\), le point \(A\); pour \(k=1\), la médiatrice de \([AB]\). Pour \(k>0\), \(k\ne1\), mettre au carré est équivalent; le lieu est un cercle d’Apollonius. Son centre est le barycentre de \((A,1),(B,-k^2)\) et son rayon vaut \(kAB/|1-k^2|\).
Pour \(A(0;0),B(3;0),k=2\), \(x^2+y^2=4((x-3)^2+y^2)\) donne \((x-4)^2+y^2=4\) : centre \((4;0)\) et rayon \(2\).
\end{frame}
\begin{frame}{Reconnaître les lieux}
Calculer la somme des poids et le barycentre si elle est non nulle. Réduire ensuite l’expression à \(sMG^2+c\). Résoudre l’équation pour \(MG^2\); une valeur positive donne un cercle, zéro un point, négative l’ensemble vide. Si \(s<0\), le sens d’une inégalité serait renversé en divisant; ne pas supposer que l’expression est une somme positive.
\end{frame}
\begin{frame}{Reconnaître les lieux}
Avec \(A(0;0)\), \(B(6;0)\) et poids \(2,1\), le barycentre est \(G(2;0)\). La réduction donne
\[
2MA^2+MB^2=3MG^2+2GA^2+GB^2=3MG^2+24.
\]
Le minimum est \(24\), atteint uniquement en \(G\). Le niveau \(12\) est vide; le niveau \(24\) est \(\{G\}\); le niveau \(51\) est le cercle de centre \(G\) et de rayon \(3\). Il faut calculer la constante de réduction, pas seulement le centre.
\end{frame}
\begin{frame}{Reconnaître les lieux}
Pour trois points, chercher un groupe dont le barycentre est facile à construire. Si ses poids ont une somme non nulle, remplacer ce groupe par son barycentre affecté de cette somme. Le barycentre final se construit alors sur une droite entre deux points. Vérifier les deux sommes avant les divisions.
\end{frame}
\begin{frame}{Reconnaître les lieux}
Quand les poids totalisent zéro, la formule du barycentre ne s’applique pas. L’expression initiale peut néanmoins se simplifier directement par Chasles ou en coordonnées. Son lieu n’est pas automatiquement vide; il faut étudier cette expression, sans inventer un centre.
\end{frame}
\end{document}
